\scanpage{73}
\[
H^{\unsure{1|1'}}(X,\mu_N)\simeq {}_N\operatorname{Pic}(X)
\]
\[
0\to {}_N\operatorname{Pic}^0(X)\to H^1(X,\mu_N)\to {}_N\mathrm{NS}(X)\to 0
\]
\margin{\tikz[baseline=-.4ex]{\draw[->] (.22,-.2) -- (0,-.2) -- (0,.85);\draw (0,.3) -- (.28,.3);} $X$ propre sur corps $k$ sép. clos, $N,l$ premiers à car. de $k$}

\[
H^1(X,\mathbb Z_l(1))\xleftarrow{\sim}
\deleted{H^1(X,A_X)}T_l(\underline{\operatorname{Pic}}{}^0_{X/k})\deleted{\unsure{H^1|B}}
\]
[Utiliser le fait que $\operatorname{Pic}^{\tau}(X)/\operatorname{Pic}^0(X)\deleted{\underline{\operatorname{Pic}}_{X/k}}$ est \uline{fini}]

En particulier, si $\underline{\operatorname{Pic}}{}^0_{X/k}$ propre, p.ex. $X$ normal, on trouve, introduisant $B_X=\underline{\operatorname{Pic}}{}^{00}_{X/k}$ et $A_X=B_X^*$, on a
\[
H^1(X,\mathbb Z_l\deleted{\unsure{(1)}})\xleftarrow{\varphi_X^*}H^1(A_X,\mathbb Z_l).
\]

\deleted{$H^{2n-1}$} Supposons $X$ lisse, purement $n$-dim. alors
\[
H^{2n-1}(X,\mathbb Z/N\mathbb Z)\simeq H^1(X,\mathbb Z/N\mathbb Z)^\vee(-n)
\]
d'où
\diagram{\begin{tikzcd}[ampersand replacement=\&,column sep=small]
0\arrow[r]\& \deleted{\unsure{B|P}}({}_N\mathrm{NS}(X))^\vee\deleted{\unsure{(-n)|(+n)}}\arrow[r]\&
 H^{2n-1}(X,\mathbb Z/N\mathbb Z)\arrow[r]\& {}_N A_X(k)(-n)\arrow[r]\arrow[d,dash,"\sim"']\&0\\
\&\&\& H^1(B_X,\mathbb Z/N\mathbb Z)\deleted{\unsure{(-n)|(+n)}}\inserted{(1-n)}\&
\end{tikzcd}}

Passant à la limite sur $l^{\unsure{\nu|n}}$
\diagram{\begin{tikzcd}[ampersand replacement=\&,column sep=small,row sep=small]
0\arrow[r]\&( \mathrm{NS}(X)[l])^\vee\deleted{\unsure{(-n)|(n)}}\arrow[r]\arrow[d,"\phantom{l}"]\&
H^{2n-1}(X,\mathbb Z_l)(n-1)\arrow[r,"\text{hom algébrique}"']\&
H^1(B_X,\mathbb Z_l)\deleted{\unsure{(1-n)|(-n)}}\arrow[r,dash]\&0\\
\&\substack{\text{comp.}\\ l\text{-primaire}}\&\&\&
\end{tikzcd}}
% archive sheet number: 72

